\section{Column Rank and Row Rank}

\begin{theorem}[Equality of Column and Row Rank]
% thm:17.c.1
\label{thm:rank_row_column_equality}
Let $A \in \M_{m \times n}(K)$. Then, $\rank_{\text{col}}(A) = \rank_{\text{row}}(A)$.
\end{theorem}

\begin{remark}[Notation and Properties of Rank]
\begin{enumerate}
    \item Because of this theorem, we will from now on denote $\rank(A) := \rank_{\text{col}}(A) = \rank_{\text{row}}(A)$.
    \item Suppose $A'$ is a row-reduced echelon matrix that is row-equivalent to $A$. Then, we have:
    \begin{equation*}
        \rank(A) = \rank_{\text{row}}(A) = \rank_{\text{row}}(A') = \text{number of pivots in } A'
    \end{equation*}
    \item Recall that for a linear map $T: V \to W$, we defined $\rank(T) := \dim \operatorname{Im}(T)$. Moreover, if $A \in \M_{m \times n}(K)$, then $\rank(T_A) = \dim \operatorname{Im}(T_A) = \rank_{\text{col}}(A)$.
\end{enumerate}
\end{remark}

\subsection{Preparation for the Proof}

\begin{lemma}[Rank Invariance under Composition with Isomorphisms]
% lemma:17.c.2
\label{lemma:rank_invariance_isomorphisms}
Let $T: V \to W$ be a linear map, and let $S: W \to W'$ and $L: P \to V$ be isomorphisms. Then:
\begin{enumerate}
    \item $\rank(S \circ T) = \rank(T)$.
    \item $\rank(T \circ L) = \rank(T)$.
\end{enumerate}
\end{lemma}

\begin{proof}
\textbf{Proof of (a):} By definition, $\rank(S \circ T) = \dim (S(T(V)))$. Since $S$ is an isomorphism, its restriction to the subspace $T(V) \subseteq W$ is injective. An injective linear map preserves the dimension of a vector space, and therefore:
\begin{equation*}
    \dim (S(T(V))) = \dim (T(V)) = \rank(T)
\end{equation*}

\textbf{Proof of (b):} By definition, $\rank(T \circ L) = \dim (T(L(P)))$. Since $L$ is an isomorphism, it is surjective, which implies that $L(P) = V$. Consequently:
\begin{equation*}
    \dim (T(L(P))) = \dim (T(V)) = \rank(T)
\end{equation*}
This completes the proof.
\end{proof}

\begin{lemma}[Invertibility of the Transpose]
% lemma:17.c.3
\label{lemma:transpose_invertibility}
Let $A \in \M_{n \times n}(K)$. Then, $A \in \GL_n(K)$ if and only if $A^T \in \GL_n(K)$.
\end{lemma}

\begin{proof}
We know that $A$ is invertible if and only if the associated linear map $T_A: K^n \to K^n$ is an isomorphism. This is true if and only if $\rank(T_A) = n$, which is equivalent to $\rank_{\text{col}}(A) = n$.

We also know that $A$ is row-equivalent to some row-reduced echelon matrix $A'$. Since row-equivalence corresponds to multiplying from the left by invertible matrices, we have that $A \in GL_n(K)$ if and only if $A' \in GL_n(K)$. Because $A'$ is an $n \times n$ row-reduced echelon matrix, $A'$ is invertible if and only if it has no zero rows. This holds if and only if $A'$ has $n$ pivots, which is equivalent to $\rank_{\text{row}}(A') = n$. Since row operations preserve the row space, this means $\rank_{\text{row}}(A) = n$.

In summary, we have established that $A \in \GL_n(K)$ if and only if $\rank_{\text{row}}(A) = n$. But by definition, the row rank of $A$ is exactly the column rank of its transpose, so $\rank_{\text{row}}(A) = \rank_{\text{col}}(A^T)$. Therefore, $A \in \GL_n(K)$ if and only if $\rank_{\text{col}}(A^T) = n$, which means $A^T$ is invertible.
\end{proof}

\begin{lemma}[Rank Invariance under Matrix Equivalence]
% lemma:17.c.4
\label{lemma:rank_invariance_equivalence}
Let $A, B \in \M_{m \times n}(K)$ such that $B = P A Q$, where $P \in \GL_m(K)$ and $Q \in \GL_n(K)$. Then:
\begin{enumerate}
    \item $\rank_{\text{col}}(A) = \rank_{\text{col}}(B)$.
    \item $\rank_{\text{row}}(A) = \rank_{\text{row}}(B)$.
\end{enumerate}
\end{lemma}

\begin{proof}
\textbf{Proof of (a):} Transitioning to the associated linear maps, we have $T_B = T_P \circ T_A \circ T_Q$. Since $P$ and $Q$ are invertible matrices, both $T_P$ and $T_Q$ are isomorphisms. Applying Lemma \ref{lemma:rank_invariance_isomorphisms}, we deduce that $\rank(T_B) = \rank(T_A)$. Recalling that the rank of a linear map given by a matrix matches the column rank of that matrix, we obtain $\rank_{\text{col}}(B) = \rank_{\text{col}}(A)$.

\textbf{Proof of (b):} Taking the transpose of the equation $B = PAQ$ yields $B^T = Q^T A^T P^T$. Since $P$ and $Q$ are invertible matrices, their transposes $P^T$ and $Q^T$ are also invertible (we have previously seen that $(P^T)^{-1} = (P^{-1})^T$). We can now apply part (a) of this Lemma, which we have already proven, to the matrices $A^T$ and $B^T$. This gives us $\rank_{\text{col}}(B^T) = \rank_{\text{col}}(A^T)$. By the definition of the transpose, this immediately translates to $\rank_{\text{row}}(B) = \rank_{\text{row}}(A)$.
\end{proof}

\subsection{Proof of the Main Theorem}

We are now ready for the proof of Theorem \ref{thm:rank_row_column_equality}.

\begin{proof}[Proof of Theorem \ref{thm:rank_row_column_equality}]
Let $r := \rank_{\text{col}}(A)$. By a previous proposition regarding the equivalence of matrices, we know that there exist invertible matrices $P \in GL_m(K)$ and $Q \in GL_n(K)$ such that:
\begin{equation*}
    P A Q = \begin{pmatrix} I_r & O \\ O & O \end{pmatrix} =: B
\end{equation*}
By Lemma \ref{lemma:17.c.4}, we know that $\rank_{\text{row}}(A) = \rank_{\text{row}}(B)$. Thus, it suffices to determine the row rank of $B$.

The matrix $B$ is of the block form $\left(\begin{smallmatrix} I_r & O \\ O & O \end{smallmatrix}\right)$. The non-zero rows of $B$ are precisely the first $r$ standard basis vectors $e_1^T, e_2^T, \dots, e_r^T \in K_{\text{row}}^n$. Since these vectors form a linearly independent set, the dimension of the row space of $B$ is exactly $r$. Consequently, $\rank_{\text{row}}(B) = r$.

Therefore, we conclude that $\rank_{\text{row}}(A) = r = \rank_{\text{col}}(A)$, completing the proof.
\end{proof}

\begin{importantremark}[Column Space versus Row Space]
Although $\rank_{\text{col}}(A) = \rank_{\text{row}}(A)$, the column space and the row space are, in general, different vector spaces. First of all, if $A \in \M_{m \times n}(K)$, then the column space is a subspace of $K^m$, whereas the row space is a subspace of $K_{\text{row}}^n$ (which is isomorphic to $K^n$). But even in the case where $m = n$ (i.e., $A \in \M_{n \times n}(K)$), these two spaces are generally distinct! They share the exact same dimension, but they are not the same space.
\end{importantremark}